% Replacement text for manuscript Section 4.4.
% Integration notes:
%   1. Update the existing qualitative Figure 7 label to
%      `fig:mixed-traffic-qualitative`, or change the references below.
%   2. Numerical macros below are frozen to the complete v20 bank
%      (bank ID converging_bank_2efd0ce2a79a425b; 40/40 valid runs).
%   4. Remove the submitted OA-CMPC panel (old Fig. 8) and the OA-CMPC row in
%      the submitted runtime table. The audited adapter is not the published
%      dual-branch method.
%   5. The old REI/jerk/CT_v aggregate plot may be retained only as a secondary
%      mechanism/trade-off diagnostic (preferably in the supplement); it is not
%      a substitute for the physical outcomes reported here.

\providecommand{\CarlaFigureSixFile}{figure6_carla_speed_clearance.pdf}

% Text-level findings from the complete v20 scene bank.
\providecommand{\CarlaEmptyShadowFinding}{All four controllers completed the
empty-shadow scenes without a collision or near collision.  DREAM reduced the
episode-mean ego speed from 31.34 to 30.12~m\,s$^{-1}$ relative to IDEAM,
corresponding to $CT_v=1.22$~m\,s$^{-1}$ [1.13, 1.32].  Its mean episode-minimum
clearance was 1.16~m [1.11, 1.21], compared with 1.18~m [1.16, 1.21] for IDEAM;
the paired difference was $-0.023$~m [$-0.071$, 0.002].  DREAM's mean peak
deceleration was $-2.51$~m\,s$^{-2}$ [$-2.58$, $-2.45$], and no empty-shadow
scene crossed the $-3$~m\,s$^{-2}$ hard-braking threshold.}
\providecommand{\CarlaTrueThreatFinding}{DREAM produced 0/5 collisions and 2/5
near collisions, compared with 4/5 and 5/5 for IDEAM, 0/5 and 5/5 for the
ADA-sourced control, and 0/5 and 4/5 for the APF-sourced control.  DREAM's mean
episode-minimum signed clearance was 1.07~m [0.94, 1.19].  The paired clearance
increases relative to IDEAM, ADA-sourced, and APF-sourced were 1.19~m
[0.89, 1.56], 0.67~m [0.52, 0.82], and 0.19~m [0.02, 0.39], respectively.  The
hidden-vehicle-specific clearance and two-dimensional TTC increased by 1.49~m
[1.06, 1.99] and 1.07~s [0.97, 1.19] relative to IDEAM.}
\providecommand{\CarlaFollowerFinding}{Relative to IDEAM, DREAM changed the
maximum follower speed loss by only $5.7\times10^{-6}$~m\,s$^{-1}$
[$-1.23\times10^{-3}$, $1.25\times10^{-3}$] and did not change the hard-braking
actor count.  The total integrated traffic speed deficit increased by
0.203~vehicle-m [0.139, 0.267], or 0.33\% of the IDEAM total.  Thus, the tested
empty shadows showed a measurable ego-level efficiency cost, but no increase
was observed in either the hard-braking count or the maximum-follower-speed-loss
metric; the 6-s horizon does not exclude longer-range traffic-wave
amplification.}

% True-threat condition.
\providecommand{\CarlaTTCollisionDREAM}{0/5}
\providecommand{\CarlaTTCollisionIDEAM}{4/5}
\providecommand{\CarlaTTCollisionADA}{0/5}
\providecommand{\CarlaTTCollisionAPF}{0/5}
\providecommand{\CarlaTTNearDREAM}{2/5}
\providecommand{\CarlaTTNearIDEAM}{5/5}
\providecommand{\CarlaTTNearADA}{5/5}
\providecommand{\CarlaTTNearAPF}{4/5}
\providecommand{\CarlaTTClearDREAM}{1.07 [0.94, 1.19]}
\providecommand{\CarlaTTClearIDEAM}{$-0.12$ [$-0.49$, 0.22]}
\providecommand{\CarlaTTClearADA}{0.41 [0.20, 0.61]}
\providecommand{\CarlaTTClearAPF}{0.89 [0.70, 1.07]}

% Empty-shadow condition.
\providecommand{\CarlaESSpeedDREAM}{30.12 [29.90, 30.34]}
\providecommand{\CarlaESSpeedIDEAM}{31.34 [31.20, 31.49]}
\providecommand{\CarlaESSpeedADA}{31.06 [30.83, 31.29]}
\providecommand{\CarlaESSpeedAPF}{31.16 [30.94, 31.37]}
\providecommand{\CarlaESCTVdream}{1.22 [1.13, 1.32]}
\providecommand{\CarlaESCTVada}{0.29 [0.19, 0.40]}
\providecommand{\CarlaESCTVapf}{0.19 [0.07, 0.33]}
\providecommand{\CarlaESClearDREAM}{1.16 [1.11, 1.21]}
\providecommand{\CarlaESClearIDEAM}{1.18 [1.16, 1.21]}
\providecommand{\CarlaESClearADA}{1.16 [1.11, 1.21]}
\providecommand{\CarlaESClearAPF}{1.17 [1.12, 1.21]}
\providecommand{\CarlaESFollowLossDREAM}{2.013 [1.626, 2.381]}
\providecommand{\CarlaESFollowLossIDEAM}{2.013 [1.627, 2.380]}
\providecommand{\CarlaESFollowLossADA}{2.014 [1.627, 2.381]}
\providecommand{\CarlaESFollowLossAPF}{2.014 [1.627, 2.382]}
\providecommand{\CarlaESHardBrakeDREAM}{0.20 [0.00, 0.60]}
\providecommand{\CarlaESHardBrakeIDEAM}{0.20 [0.00, 0.60]}
\providecommand{\CarlaESHardBrakeADA}{0.20 [0.00, 0.60]}
\providecommand{\CarlaESHardBrakeAPF}{0.20 [0.00, 0.60]}

% Runtime values from the five DREAM true-threat arms in the v20 bank.
\providecommand{\CarlaHighLevelMeanMs}{442.1}
\providecommand{\CarlaHighLevelPNinetyFiveMs}{521.0}
\providecommand{\CarlaHighLevelMaxMs}{536.1}
\providecommand{\CarlaThreatMeanSpeed}{29.92}
\providecommand{\CarlaTravelAtMeanLatency}{13.2}
\providecommand{\CarlaTravelAtPNinetyFiveLatency}{15.6}
\providecommand{\CarlaTravelAtMaxLatency}{16.0}
\providecommand{\CarlaHighLevelRateHz}{2.33}
\providecommand{\CarlaDroppedRequestsPct}{76.3}
\providecommand{\CarlaLowLevelMeanMs}{1.493}
\providecommand{\CarlaLowLevelPNinetyFiveMs}{2.029}
\providecommand{\CarlaLowLevelMaxMs}{3.420}
\providecommand{\CarlaPhysicsMeanMs}{12.557}
\providecommand{\CarlaPhysicsPNinetyFiveMs}{16.574}
\providecommand{\CarlaPhysicsMaxMs}{20.677}
\providecommand{\CarlaRevealPlanDelayMeanS}{0.85}
\providecommand{\CarlaRevealPlanDelayLowS}{0.80}
\providecommand{\CarlaRevealPlanDelayHighS}{0.95}
\providecommand{\CarlaRevealPlanDelayMaxS}{1.05}
\providecommand{\CarlaRevealPlanTravelMeanM}{25.4}
\providecommand{\CarlaRevealPlanTravelMaxM}{31.4}

\subsection{Evaluation results}
\label{sec:evaluation-results}

\subsubsection{Mechanism-level qualitative evidence}
\label{sec:mechanism-qualitative}

Figure~\ref{fig:mixed-traffic-qualitative} examines how the alternative field
sources represent an occluded conflict in mixed traffic.  The snapshots show
the instantaneous spatial support of each field and its relation to the
corresponding ego response.  DREAM places caution inside the shadow before a
hidden actor is observed.  The ADA-sourced and APF-sourced controls produce
different spatial patterns under the same shared planning backbone.  These
examples expose the mechanism and support visual comparison of the field
geometry.  They do not estimate event incidence, establish a safety margin,
or provide a guarantee for unobserved traffic scenes.

The qualitative visualization and the closed-loop experiment serve distinct
roles.  Figure~\ref{fig:mixed-traffic-qualitative} makes the field mechanism
visible in a representative mixed-traffic interaction.  Figure~\ref{fig:carla-two-condition-profiles}
instead aggregates matched, high-fidelity CARLA runs.  It tests whether the
mechanism changes physical responses when the same shadow is empty or contains
a conflicting vehicle.

\subsubsection{Matched empty-shadow and true-threat comparison}
\label{sec:carla-two-condition}

The CARLA experiment used a converging overtaking scenario in Town06.  The ego
vehicle passed CARLA's rigid fire-engine asset, which served as a large truck
surrogate rather than an articulated truck--trailer.  A second vehicle passed
on the opposite side.  Both vehicles then requested entry into the centre lane
ahead of the occluder.  Surrounding vehicles followed an
intelligent-driver-model controller.  The \emph{empty-shadow} condition removed
the hidden vehicle, whereas the \emph{true-threat} condition retained its
conflicting maneuver.

Five randomized scene blocks were evaluated with DREAM, IDEAM, an ADA-sourced
field control, and an APF-sourced field control.  The resulting bank contained
$5\times2\times4=40$ runs.  Within each block, all eight arms used the same
frozen nominal manifest, random seed, actor definitions, and background-traffic
control laws.  Realized
initial states were admitted only when they satisfied the common tolerances,
and construction hashes were required to match across arms.  Arm order was
randomized within each block.  The ADA-sourced and APF-sourced arms changed
the field source while retaining the common propagation, coupling, and
MPC--CBF backbone.  They are
therefore source-level controls rather than complete reimplementations of the
cited external planners.  OA-CMPC was not included because the available
adapter does not reproduce its original dual-branch contingency optimization.

The true-threat traces were aligned to the first semantic-LiDAR detection of
the hidden vehicle.  The matching empty-shadow trace used the reveal time from
the corresponding true-threat arm.  This construction preserves a common
event origin without introducing a synthetic reveal into the empty condition.
The plotted centre line is the mean over complete scene blocks.  Shading shows
a descriptive 95\% scene-block bootstrap interval from 10,000 resamples.  The
scene block, rather than the 20-Hz simulation tick, is the resampling unit.
The intervals summarize this finite scenario bank and are not used for
population-level hypothesis testing.

For each neighboring vehicle $i$, let $B_{\mathrm{e}}(t)$ and $B_i(t)$ denote
their ground-plane oriented bounding boxes (OBBs).  Their signed separation is
\begin{equation}
\delta_i(t)=
\begin{cases}
\displaystyle \min_{x\in B_{\mathrm{e}}(t),\,y\in B_i(t)}\lVert x-y\rVert_2,
& B_{\mathrm{e}}(t)\cap B_i(t)=\varnothing,\\[2mm]
-\operatorname{mtd}\!\left(B_{\mathrm{e}}(t),B_i(t)\right),
& B_{\mathrm{e}}(t)\cap B_i(t)\neq\varnothing,
\end{cases}
\qquad
d_{\min}(t)=\min_i\delta_i(t),
\label{eq:signed-obb-clearance}
\end{equation}
where $\operatorname{mtd}$ is the minimum translation distance required to
separate two intersecting boxes.  Positive values indicate separation, zero
indicates contact, and negative values indicate footprint overlap.  This
definition avoids the ambiguity of centre-to-centre distance for vehicles
with different dimensions.

Collision incidence is obtained from the CARLA collision sensor.  A
near-collision is recorded when the positive signed clearance falls below
0.5~m or the constant-velocity oriented-box TTC falls below 1.0~s.  The two
indicators are reported independently; an episode that later collides may
also contain a preceding near-collision interval.

Figure~\ref{fig:carla-two-condition-profiles}(a,b) reports the empty-shadow
response.  \CarlaEmptyShadowFinding  The empty-shadow conservatism tax for
controller $m$ is computed within each scene block as
$CT_v^{(m)}=\bar v_{\mathrm{ego}}^{\mathrm{IDEAM}}-
\bar v_{\mathrm{ego}}^{(m)}$.  It therefore quantifies defensive speed loss
relative to the no-field shared-backbone reference.  The follower-disturbance
metrics in Table~\ref{tab:carla-two-condition} assess short-horizon traffic
disturbance rather than long-range traffic-wave formation.  \CarlaFollowerFinding

Figure~\ref{fig:carla-two-condition-profiles}(c,d) reports the corresponding
true-threat response.  \CarlaTrueThreatFinding  Collision and near-collision
incidence are reported separately from the continuous clearance trace.  This
separation prevents a smooth average from masking discrete adverse events.

\begin{figure*}[t]
\centering
\includegraphics[width=\textwidth]{\CarlaFigureSixFile}
\caption{Matched closed-loop response to empty and occupied occlusion shadows.
(a,b) Ego speed and signed nearest-neighbor oriented-box clearance in the
empty-shadow condition.  (c,d) The corresponding profiles in the true-threat
condition.  Faint curves show individual scene blocks, bold curves show their
means, and bands show descriptive 95\% scene-block bootstrap intervals.  Time
zero marks semantic-LiDAR reveal in the true-threat arm and its matched event
time in the empty-shadow arm.}
\label{fig:carla-two-condition-profiles}
\end{figure*}

\begin{table*}[t]
\centering
\caption{Matched two-condition CARLA results over five randomized scene
blocks.  Incidence entries report episodes out of five.  Continuous entries
report the scene-block mean [descriptive 95\% bootstrap interval].  No
population-level significance inference is made.}
\label{tab:carla-two-condition}
\scriptsize
\setlength{\tabcolsep}{3pt}
\resizebox{\textwidth}{!}{%
\begin{tabular}{llcccc}
\toprule
Condition & Outcome & DREAM & IDEAM & ADA-sourced & APF-sourced \\
\midrule
True threat
& Collision incidence
& \CarlaTTCollisionDREAM & \CarlaTTCollisionIDEAM
& \CarlaTTCollisionADA & \CarlaTTCollisionAPF \\
& Near-collision incidence
& \CarlaTTNearDREAM & \CarlaTTNearIDEAM
& \CarlaTTNearADA & \CarlaTTNearAPF \\
& Minimum signed OBB clearance (m)
& \CarlaTTClearDREAM & \CarlaTTClearIDEAM
& \CarlaTTClearADA & \CarlaTTClearAPF \\
\midrule
Empty shadow
& Mean ego speed (m\,s$^{-1}$)
& \CarlaESSpeedDREAM & \CarlaESSpeedIDEAM
& \CarlaESSpeedADA & \CarlaESSpeedAPF \\
& $CT_v$ relative to IDEAM (m\,s$^{-1}$)
& \CarlaESCTVdream & 0
& \CarlaESCTVada & \CarlaESCTVapf \\
& Minimum signed OBB clearance (m)
& \CarlaESClearDREAM & \CarlaESClearIDEAM
& \CarlaESClearADA & \CarlaESClearAPF \\
& Maximum follower speed loss (m\,s$^{-1}$)
& \CarlaESFollowLossDREAM & \CarlaESFollowLossIDEAM
& \CarlaESFollowLossADA & \CarlaESFollowLossAPF \\
& Hard-braking traffic actors (count)
& \CarlaESHardBrakeDREAM & \CarlaESHardBrakeIDEAM
& \CarlaESHardBrakeADA & \CarlaESHardBrakeAPF \\
\bottomrule
\end{tabular}
}
\end{table*}

\subsubsection{Asynchronous execution and latency exposure}
\label{sec:asynchronous-execution}

The current high-level optimizer remained slower than its nominal 100-ms
request period.  Its mean, mean per-run 95th percentile, and maximum service
times were \CarlaHighLevelMeanMs, \CarlaHighLevelPNinetyFiveMs, and
\CarlaHighLevelMaxMs~ms, respectively.  At the mean true-threat ego speed of
\CarlaThreatMeanSpeed~m\,s$^{-1}$, these durations corresponded to
\CarlaTravelAtMeanLatency, \CarlaTravelAtPNinetyFiveLatency, and
\CarlaTravelAtMaxLatency~m of travel.  The scheduler coalesced pending
requests and delivered \CarlaHighLevelRateHz{} completed updates per second.
It discarded \CarlaDroppedRequestsPct\% of requests that had been superseded
before execution.

The low-level controller ran at 10~Hz while tracking the latest time-stamped
plan.  Its mean, 95th-percentile, and maximum execution times were
\CarlaLowLevelMeanMs, \CarlaLowLevelPNinetyFiveMs, and \CarlaLowLevelMaxMs~ms.  The
20-Hz physics and control loop required \CarlaPhysicsMeanMs,
\CarlaPhysicsPNinetyFiveMs, and \CarlaPhysicsMaxMs~ms, respectively.  Neither loop
missed its scheduled deadline in this experiment.  Lane holding and
visible-actor emergency braking remained available while a new plan was
pending.  This tracker and fallback layer was identical in all CARLA arms;
the controller comparison changed the high-level plan supplied to that common
execution layer.

\begin{table*}[t]
\centering
\caption{Measured asynchronous execution in the five DREAM true-threat runs.
The 95th-percentile column is the mean of the five within-run percentiles, and
the maximum column is the largest observed value across all five runs.  Timing
runs were real-time paced and excluded frame encoding.}
\label{tab:carla-asynchronous-runtime}
\small
\begin{tabular}{lccccc}
\toprule
Layer & Nominal period & Mean (ms) & Mean P95 (ms) & Maximum (ms) & Scheduling outcome \\
\midrule
High-level optimizer & 100 ms request & 442.1 & 521.0 & 536.1 & 2.33 Hz; 76.3\% coalesced \\
Low-level controller & 100 ms & 1.493 & 2.029 & 3.420 & 0/300 deadline misses \\
Physics/control loop & 50 ms & 12.557 & 16.574 & 20.677 & 0/600 deadline misses \\
\bottomrule
\end{tabular}
\end{table*}

Timely low-level execution limits actuation delay but does not remove
high-level planning latency.  A hidden actor can reveal while the controller
is still tracking a plan computed from older information.  The measured mean
delay from reveal to application of a hidden-aware DREAM plan was
\CarlaRevealPlanDelayMeanS~s
with a descriptive 95\% scene-block bootstrap interval of
[\CarlaRevealPlanDelayLowS, \CarlaRevealPlanDelayHighS]; the largest observed
delay was \CarlaRevealPlanDelayMaxS~s.  At the reported mean ego
speed, these end-to-end delays correspond to approximately
\CarlaRevealPlanTravelMeanM{} and \CarlaRevealPlanTravelMaxM{}~m of travel,
respectively.  The 10-Hz layer must
therefore be interpreted as an execution mechanism, not as a safety
certificate.  These measurements demonstrate an implemented asynchronous
proof of concept.  They do not establish 10-Hz feasibility of the current
high-level optimizer or guaranteed safety during an aggressive reveal.
